\section{MLIR 与 AscendNPU IR 静态探索}
\label{sec:mlir}

LLVM IR 适合表达接近通用机器模型的程序，但领域信息若过早压平，后续优化难以恢复。多级中间表示（Multi-Level Intermediate Representation, MLIR）允许方言（dialect）、操作、类型和属性在多个抽象层共存，并用方言转换的合法化目标和类型转换逐步降低\cite{lattner2021mlir,mlirDialectConversion}。渐进式降低（progressive lowering）的工程价值不是自动保证正确或快速，而是把每次语义损失变成显式、可检查的边界。

进阶目录中的 \texttt{add.mlir} 保存 AscendNPU IR 官方 VecAdd 数据通路的最小静态资产\cite{ascendVecAdd}。结构检查确认以下顺序与属性同时存在：
\begin{enumerate}
  \item \texttt{hacc.entry} 与 \texttt{DEVICE} 标记设备 kernel 入口；
  \item 两个 \texttt{hivm.hir.load} 把 16 个 \texttt{i16} 元素从全局内存（\texttt{gm}）搬到统一缓冲区（\texttt{ub}）；
  \item \texttt{hivm.hir.vadd} 在 \texttt{ub} 层执行向量加法；
  \item \texttt{hivm.hir.store} 把结果从 \texttt{ub} 写回 \texttt{gm}。
\end{enumerate}
该表示在降到 LLVM IR 前仍显式保留设备入口、存储层次、搬运与向量操作，展示了多级 IR 相比立即压成指针算术的优势。

必须区分“该快照内部的数据通路”和“多个快照之间的降低”。官方架构给出的文档链路为：HFusion 保留 named-operation 高层语义并完成预处理/融合；转换到 HIVM 后逐步落实核间映射、片上内存与处理单元/同步；低层 MLIR 再由 \texttt{hivmc} 转成 LLVM IR 并生成算子二进制\cite{ascendArchitecture}。本实验保存的 \texttt{add.mlir} 已处于 HIVM 层，故 \texttt{gm $\rightarrow$ ub $\rightarrow$ vadd $\rightarrow$ gm} 是同一层的领域数据流，不是本地观测到的 progressive-lowering 轨迹。表~\ref{tab:ascend-lowering} 将官方说明与本地证据明确分开。

\begin{table}
\caption{AscendNPU IR 降低边界及本实验的证据强度}
\label{tab:ascend-lowering}
\centering
\begin{tabularx}{\textwidth}{lXX}
\toprule
边界 & 被具体化的语义 & 本地状态 \\
\midrule
HFusion $\rightarrow$ HIVM & 从高层 named operation 到 NPU 计算、搬运与同步 & 官方文档；未执行 \\
HIVM 优化 & 核映射、片上内存、处理单元与同步 & 官方文档；仅保存输入快照 \\
低层 MLIR $\rightarrow$ LLVM IR & 通用后端可消费的低层表示 & 官方文档；未执行 \\
LLVM IR $\rightarrow$ 算子二进制 & 目标指令选择与设备对象生成 & 官方文档；未执行 \\
\bottomrule
\end{tabularx}
\end{table}

\texttt{advanced\_structure} 测试检查关键方言标记以及 \texttt{load $\rightarrow$ vadd $\rightarrow$ store} 的文本顺序，并在 Linux CTest 中通过。其证据强度仅为“教学资产包含预期静态结构”：普通 \texttt{mlir-opt} 不认识 \texttt{hacc}/\texttt{hivm} 方言；完整路径还需要 CANN 环境、\texttt{bishengir-compile}、匹配的驱动和 Ascend NPU。本文未执行方言解析、专用降低、生成设备对象或上板验证，因此不声称 VecAdd 已编译或运行。
